\chapter*{Resumen}

RISC-V es una especificación de conjunto de instrucciones o ISA surgida en el año 2010 en la Universidad de Berkeley como una alternativa libre a los conjuntos de instrucciones (privativos) dominantes en aquel momento. Nacida de un proyecto académico destinado a facilitar la investigación en arquitecturas hardare, su modelo abierto y estandarizado, así como su filosofía, han favorecido su consolidación como una alternativa real a arquitecturas tradicionales como x86 o ARM en determinadas áreas.

Este documento pretende recoger el proceso de implementación de un núcleo RISC-V segmentado en cinco etapas, capaz de ejecutar el conjunto de instrucciones base definido en el estándar.

El proyecto tiene como punto de partida un diseño monociclo incompleto, desarrollado en la Universidad Nacional Cheng Kung (Taiwán). Esta implementación ha sido completada y estudiada en profundidad, quedando plasmadas en este documento las conclusiones de ese proceso de indagación.

Por último, se detalla el proceso de despliegue del diseño final sobre el chip FPGA de una placa de desarrollo usando el software Vivado.\newline\newline

\textbf{Palabras clave:} RISC-V, Chisel, HDL, arquitectura de computadores, diseño hardware.